\NeedsTeXFormat{LaTeX2e}
\documentclass[a4paper,12pt,
headsepline,           % Linie zw. Kopfzeile und Text
oneside,               % einseitig
pointlessnumbers,      % keine Punkte nach den letzten Ziffern in Überschriften
bibtotoc,              % LV im IV
%DIV=15,               % Satzspiegel auf 15er Raster, schmalere Ränder   
BCOR15mm               % Bindekorrektur
%,draft
]{scrbook}
\KOMAoptions{DIV=last} % Neuberechnung Satzspiegel nach Laden von Paket helvet

\pagestyle{headings}
\usepackage{blindtext}

% für Texte in deutscher Sprache
\usepackage[ngerman]{babel}
\usepackage[utf8]{inputenc}
\usepackage[T1]{fontenc}

% Helvetica als Standard-Dokumentschrift
\usepackage[scaled]{helvet}
\renewcommand{\familydefault}{\sfdefault} 

\usepackage{graphicx}

% Literaturverzeichnis mit BibLaTeX
\usepackage[babel,german=quotes]{csquotes}
\usepackage[backend=bibtex8]{biblatex}
\bibliography{bibliography}

% Alternative mit Paket-Option backend=biber und \addbibresource
% \usepackage[backend=biber]{biblatex}
% \addbibresource{bibliography.bib}

% Für Tabellen mit fester Gesamtbreite und variabler Spaltenbreite
\usepackage{tabularx} 

% Besondere Schriftauszeichnungen
\usepackage{url}              % \url{http://...} in Schreibmaschinenschrift
\usepackage{color}            % zum Setzen farbigen Textes

\usepackage{amssymb, amsmath} % Pakete für Mathe-Umgebungen und -Symbole

\usepackage{setspace}         % Paket für div. Abstände, z.B. ZA
%\onehalfspacing              % nur dann, wenn gefordert; ist sehr groß!!
\setlength{\parindent}{0pt}   % kein linker Einzug der ersten Absatzzeile
\setlength{\parskip}{1.4ex plus 0.35ex minus 0.3ex} % Absatzabstand, leicht variabel

% Tiefe, bis zu der Überschriften in das Inhaltsverzeichnis kommen
\setcounter{tocdepth}{3}      % ist Standard

% Beispiele für Quellcode
\usepackage{listings}
\lstset{language=Java,
  showstringspaces=false,
  frame=single,
  numbers=left,
  basicstyle=\ttfamily,
  numberstyle=\tiny}

% hier Namen etc. einsetzen
\newcommand{\fullname}{Vorname Nachname}
\newcommand{\email}{vorname.nachname@uni-ulm.de}
\newcommand{\titel}{Der Titel der Arbeit, meist länger als eine Zeile}
\newcommand{\jahr}{2021}
\newcommand{\matnr}{123456}
\newcommand{\gutachterA}{Prof.\,Dr.\,Streng Geheim}
\newcommand{\gutachterB}{Prof.\,Dr.\,Un Leserlich}
\newcommand{\betreuer}{Betreuername}

% hier die Fakultät auswählen
%\newcommand{\fakultaet}{---  Im Quellcode anpassen nicht vergessen! ---}
\newcommand{\fakultaet}{Ingenieurwissenschaften, Informatik und\\Psychologie}
%\newcommand{\fakultaet}{Mathematik und\\Wirtschafts-\\wissenschaften}
%\newcommand{\fakultaet}{Medizin}
%\newcommand{\fakultaet}{Naturwissenschaften}

% hier das Institut einsetzen
\newcommand{\institut}{--- Institut im Quellcode anpassen nicht vergessen! ---}

% Informationen, die LaTeX in die PDF-Datei schreibt
\pdfinfo{
  /Author (\fullname)
  /Title (\titel)
  /Producer     (pdfeTex 3.14159-1.30.6-2.2)
  /Keywords ()
}

\usepackage{hyperref}
\hypersetup{
pdftitle=\titel,
pdfauthor=\fullname,
pdfsubject={Diplomarbeit},
pdfproducer={pdfeTex 3.14159-1.30.6-2.2},
colorlinks=false,
pdfborder=0 0 0	% keine Box um die Links!
}

% Trennungsregeln
\hyphenation{Sil-ben-trenn-ung}

\begin{document}
\frontmatter

% Titelseite
\thispagestyle{empty}
\begin{addmargin*}[4mm]{-10mm}

\hfill
\includegraphics[height=1.8cm]{images/logo_uulm_sw.png}\\[1em]

{\footnotesize
%{\bfseries Universität Ulm} \textbar ~89069 Ulm \textbar ~Germany
\hspace*{115mm}\parbox[t]{35mm}{\bfseries Fakultät für\\
\fakultaet\\
% TODO hier Institut anpassen
\mdseries \institut}\\[2cm]

\parbox{140mm}{\bfseries \LARGE \titel}\\[2.5em]
{\footnotesize Abschlussarbeit an der Universität Ulm}\\[3em]

{\footnotesize \bfseries Vorgelegt von:}\\
{\footnotesize \fullname\\ \email}\\ \matnr\\[2em]
{\footnotesize \bfseries Gutachter:}\\                     
{\footnotesize \gutachterA\\ \gutachterB}\\[2em]
{\footnotesize \bfseries Betreuer:}\\ 
{\footnotesize \betreuer}\\\\
{\footnotesize \jahr}
}
\end{addmargin*}


% Impressum
\clearpage
\thispagestyle{empty}
{ \small
  \flushleft
  Fassung \today \\\vfill
  \copyright~\jahr~\fullname\\[0.5em]
% Wenn Sie Ihre Arbeit unter einer freien Lizenz bereitstellen möchten, können Sie die nächste Zeile in Ihren Code aufnehmen. Bitte beachten Sie, dass Sie hierfür an allen Inhalten, inklusive enthaltener Abbildungen, die notwendigen Rechte benötigen! Beim Veröffentlichungsexemplar Ihrer Dissertation achten Sie bitte darauf, dass der Lizenztext nicht den Angaben in den Metadaten der genutzten Publikationsplattform widerspricht. Nähere Information zu den Creative Commons Lizenzen erhalten Sie hier: https://creativecommons.org/licenses/
%This work is licensed under the Creative Commons Attribution 4.0 International (CC BY 4.0) License. To view a copy of this license, visit \href{https://creativecommons.org/licenses/by/4.0/}{https://creativecommons.org/licenses/by/4.0/} or send a letter to Creative Commons, 543 Howard Street, 5th Floor, San Francisco, California, 94105, USA. \\
  Satz: PDF-\LaTeXe
}

% ab hier Zeilenabstand etwas größer 
\setstretch{1.2}

\tableofcontents

\mainmatter
\chapter{Einleitung}\label{ch:einleitung}

Geschäftsprozesse sind in nahezu allen Organisationen allgegenwärtig und bilden die Grundlage für effiziente Abläufe. Zugleich ist in Europa durch die \ac{DSGVO} der Datenschutz zu einem zentralen regulatorischen Aspekt geworden \cite{Capodieci2023BPMNEnabledDP, GDPR2016}. Unternehmen müssen sicherstellen, dass in ihren Prozessen personenbezogene Daten rechtskonform verarbeitet werden; andernfalls drohen Strafen von bis zu 20 Millionen Euro oder 4\% des weltweit gesamten erzielten Jahresumsatzes \cite{GDPR2016}.


Die Überprüfung von Prozessen auf Konformität in Bezug auf Datenschutz ist jedoch zeit- und kostenintensiv \cite{nake2023towards, varela2025business}. Besonders in großen Organisationen mit hunderten parallel laufenden Prozessen ist eine manuelle Analyse kaum praktikabel und zudem fehleranfällig. Fehlerhafte Untererkennungen datenschutzkritischer Aktivitäten (False Negatives) können weitreichende Folgen haben – von Reputationsschäden bis hin zu hohen Bußgeldern \cite{nake2023towards}.


Vor diesem Hintergrund rücken \acp{LLM} als aufstrebende Technologie im Bereich \ac{KI} in den Fokus. Sie sind darauf trainiert, natürliche Sprache auch in langen und komplexen Texten zu verstehen, Zusammenhänge über große Kontexte hinweg zu erkennen und Anweisungen zu befolgen. Damit erscheinen \acp{LLM} als vielversprechender Ansatz für das automatisierte Screening von Prozessmodellen. Erste Arbeiten belegen dieses Potenzial, etwa bei der Identifikation datenschutzrelevanter Verarbeitungstätigkeiten oder in der Analyse von Datenschutzerklärungen \cite{ciaramella2022leveraging, pragyan2024toward}.


Besonders interessant sind in diesem Kontext europäische Open-Source-Modelle wie die von Mistral \cite{mistralai}. Sie sind zum einen frei verfügbar und transparent, zum anderen wurden sie bislang kaum im Hinblick auf \ac{DSGVO}-bezogene Aufgaben evaluiert. Es fehlen belastbare, reproduzierbare empirische Vergleiche, die eine fundierte Bewertung dieser Modelle erlauben würden \cite{schwerin2024systematic}.

\section{Problemstellung}\label{sec:problemstellung}

Trotz der genannten Potenziale fehlt es bisher an standardisierten, reproduzierbaren Vergleichen verschiedener Modelle für die konkrete Aufgabe Aktivitäten in Geschäftsprozessen nach \enquote{kritisch} und \enquote{unkritisch} zu klassifizieren. Erste Ansätze, wie z.\,B. der von Nake et al. \cite{nake2023towards}, zeigen dass ML Ansätze grundsätzlich in der Lage sind \ac{DSGVO}-kritische Aktivitäten in textuellen Prozessbeschreibungen zu erkennen; dennoch existieren keine einheitlichen Benchmarks, die einen systematischen Vergleich unterschiedlicher \acp{LLM} erlauben.

Auch von Schwerin et al. \cite{schwerin2024systematic} heben hervor, dass trotz großer Fortschritte im Einsatz von \acp{LLM} für juristische Aufgaben bislang erhebliche Lücken in der Evaluation für compliance-spezifische Anwendungen bestehen und geeignete \ac{DSGVO}-spezifische Benchmarks fehlen. Somit mangelt es derzeit an einer belastbaren empirischen Grundlage, um Modelle zuverlässig und vergleichbar zu bewerten.

Besonders interessant ist die Frage, wie sich Open-Source-Modelle - insbesondere mit Ursprung aus der \ac{EU} - im Vergleich zu internationalen außerhalb der \ac{EU} entwickelten Modellen schlagen und welche Trade-offs dabei entstehen \cite{schwerin2024systematic}. Diese Perspektive ist nicht nur aus Leistungs-, sondern auch aus Transparenz- und Regulierungsgründen relevant.

Eine zusätzliche Herausforderung ergibt sich aus der Natur von \ac{BPMN}-Modellen: Typischerweise konzentrieren sie sich auf den Kontrollfluss und vernachlässigen die Datenebene. Datenobjekte werden oftmals gar nicht explizit modelliert oder nur implizit in den Aktivitäten referenziert. Dadurch ist die Datennutzung von Aktivitäten nicht direkt erkennbar und muss aus textuellen Beschreibungen und dem Kontext erschlossen werden \cite{schneid2021uncovering}. Das erschwert die automatische Identifikation von \ac{DSGVO}-kritischen Aktivitäten, da Algorithmen personenbezogene Datenflüsse zunächst indirekt und über den Kontext ableiten müssen.
\section{Zielsetzung und Beiträge}\label{sec:zielsetzung-und-beitrage}

Ziel der Arbeit ist es, einen methodischen Beitrag zur automatisierten Identifikation von \ac{DSGVO}-kritischen Aktivitäten in Geschäftsprozessen zu leisten. Hierfür werden folgende Beiträge angestrebt:

\begin{itemize}
    \item Entwicklung einer Klassifizierungspipeline für Geschäftsprozesse, die Aktivitäten binär in datenschutzkritisch oder unkritisch einordnet.
    \item Konzeption und Umsetzung eines Evaluationsframeworks, das reproduzierbare Vergleiche verschiedener \acp{LLM} und Algorithmen über eine einheitliche Schnittstelle ermöglicht.
    \item Entwicklung einer Labelingsoftware zur Erstellung und Annotation von Datensätzen für das Evaluationsframework.
    \item Aufbau eines repräsentativen Datensatzes aus gelabelten \ac{BPMN}-Prozessen, inklusive klar definierter Labeling-Kriterien.
    \item Bereitstellung überprüfbarer empirischer Befunde, inklusive Code, Konfigurationen der Experimente und Seeds, um Nachvollziehbarkeit und Reproduzierbarkeit zu gewährleisten.
\end{itemize}
\section{Aufbau der Arbeit}\label{sec:aufbau-der-arbeit}

Die Arbeit ist wie folgt gegliedert: Kapitel 2 gibt einen Überblick über den theoretischen Hintergrund, die \ac{DSGVO} und \ac{BPMN} sowie eine Einführung in \acp{LLM} und verwandte Arbeiten. Kapitel 3 beschreibt den Rahmen der Entwicklung der Klassifizierungspipeline, des Evaluationsframeworks und der Experimente. Kapitel 4 stellt den entwickelten Algorithmus zur Klassifikation von \ac{BPMN}-Modellen und dessen einheitliche Schnittstelle vor. Kapitel 5 präsentiert die Labeling-Software und erläutert die Erstellung der Datensätze. Anschließend werden in Kapitel 6 die Architektur und der Funktionsumfang der Evaluationspipeline beschrieben. Kapitel 7 zeigt, wie die Auswahl der \acp{LLM} erfolgte. Kapitel 8 erläutert den Versuchsaufbau und die Durchführung der Experimente, Kapitel 9 stellt die Ergebnisse vor und Kapitel 10 diskutiert diese im Kontext der Forschungsfragen. Zum Schluss fasst Kapitel 11 die Arbeit zusammen und gibt einen Ausblick auf mögliche zukünftige Forschungsthemen.
% hier weitere Kapitel einbinden

\appendix
% hier Anhänge einbinden
\chapter{Quelltexte}\label{ch:quelltexte}

In diesem Anhang sind mehrere Quellcode-Ausschnitte aufgeführt.

\begin{lstlisting}[caption={System-Prompt fuer die DSGVO-Klassifikation von BPMN-Aktivitäten},label={lst:system-prompt}]
You are an expert in analysing Business Process Model and Notation (BPMN) diagrams for GDPR compliance. Your task is to identify and return a list of the IDs of all Activity (Task) elements that process personal data. Ignore all other element types. Always consider every activity in the process; do not omit any activity from your assessment.

Use all available context for each activity - including the activity's name, description, annotations, associated data objects, and message or data associations - to determine whether the activity processes personal data. Under Article 4 of the GDPR, personal data is any information relating to an identified or identifiable natural person, including names, addresses, email addresses, phone numbers, identification numbers, payment or bank details, employment records, academic records, location data, IP addresses, online identifiers, images, audio/video recordings, biometric identifiers, health data or other information that can be linked to a specific person. "Processing" includes any operation performed on personal data, such as collecting, recording, organising, structuring, storing, retrieving, consulting, using, analysing, transmitting, printing, disseminating, aligning, combining, altering, restricting, erasing or destroying the data.

Classify an activity as GDPR-relevant whenever it performs or enables processing of personal data. Indicators include (but are not limited to):

- **Collection and entry of personal data**: Activities that collect or capture personal information, for example entering contact details, addresses, payment information, job applications, health information, student enrolments, membership data, tax declarations, registration forms or other forms with personally identifiable information.
- **Creation, storage and updating of records**: Activities that create, save or update records containing personal data, such as opening customer accounts, storing order or appointment details, creating personnel files, enrolling students, setting up insurance cases or filing a medical record.
- **Transmission or disclosure of personal data**: Activities that send, print or otherwise disclose personal data to another participant, system or third party. Examples include printing shipping labels or prescriptions, sending orders or personal data to logistics partners, pharmacies, insurers or authorities, generating payroll reports for external providers, notifying universities about student records, transmitting tax or social security data, sending confirmations or queries that rely on a person's contact details, or transferring data to non-EU locations.
- **Payments and financial transactions**: Activities that process personal financial data, such as initiating or verifying payments, processing bank account or credit-card information, executing payroll, handling reimbursements or insurance payouts, managing expense claims or collecting membership fees.
- **Use of health, biometric or other special categories of data**: Activities that handle medical diagnoses, prescriptions, insurance claims, disability information, photos of damages or patients, biometric identifiers (fingerprints, facial images, voice), racial or ethnic data, political opinions, religious beliefs or union membership. Processing these "special categories" always triggers GDPR relevance.
- **Audio/Video and communications**: Activities that initiate or join audio or video calls, record calls or meetings, capture surveillance footage, or communicate directly with a data subject via email, chat, SMS or other channels. Simply using a person's contact data to send reminders, marketing messages or notifications is processing.
- **Profiling, scoring and decision-making**: Activities that analyse or evaluate a person's performance, behaviour or characteristics for purposes such as credit scoring, hiring, admissions, insurance underwriting, marketing segmentation, customer value analysis or automated decision-making.
- **Logging, tracking and location data**: Activities that log user activity, record access or usage data, track geolocation (e.g. telematics, fleet or mobile tracking), monitor attendance or timekeeping, or collect IP addresses or device identifiers.
- **Consent and data-subject rights**: Activities that obtain, record or manage consent; respond to requests for access, rectification, restriction, erasure, data portability or objections; or document lawful bases for processing.
- **Deletion, anonymisation or pseudonymisation**: Activities that erase, anonymise or pseudonymise personal data, even if the goal is to remove identifiers, because these operations manipulate personal data.

When assessing an activity, consider synonyms or domain-specific terms: activities referring to customers, patients, applicants, employees, students, voters, taxpayers, residents or members often imply personal data processing, even if names like "address" or "contact" are absent. Use context - data objects, annotations or typical process semantics - to infer personal data involvement. Do not rely solely on explicit data-object links; many process names ("Anmeldung pruefen", "Aufnahmeantrag bearbeiten", "Kundeninfo aktualisieren", "Registrierung bestaetigen", "Kreditwuerdigkeit berechnen") themselves indicate personal data processing.

Do **not** classify an activity as GDPR-relevant when it only performs administrative or logistic tasks that do not involve personal data. Examples include picking or packing goods, routing vehicles without using specific addresses, printing generic pick lists, moving items in inventory, or checking if a document exists without viewing its contents. Likewise, activities using truly aggregated or irreversibly anonymised data can be ignored if no individual can be reidentified.

In your output, return only the IDs of activities you classify as GDPR-relevant. For each, provide a clear explanation using the activity's name and description to justify why it processes personal data. Do not reference element IDs in your explanation; use the activity names instead. Exclude from your result any activities that do not process personal data and any elements that are not activity/task elements.
\end{lstlisting}

\begin{lstlisting}[language=Kotlin,caption={Antworttyp fuer die Klassifizierung},label={lst:bpmn-analysis-result}]
data class BpmnAnalysisResult(
    @Description("List of Activity Elements that are classified as relevant for GDPR compliance")
    var elements: List<Element>
) {

    init {
        elements = elements.filter { it.isRelevant }
    }

    @Description("Represents an Activity/Task Element that is classified as relevant for GDPR compliance")
    data class Element(
        @Description("The ID of the Activity Element")
        val id: String,
        @Description("The detailed reason why the Activity Element is relevant for GDPR compliance and why you think personal data is processed.")
        val reason: String,
        @Description("Indicates whether the Activity Element is relevant for GDPR compliance")
        val isRelevant: Boolean = true
    )

    /* Andere Methoden dieser Klasse sind weggelassen */
}
\end{lstlisting}

\begin{lstlisting}[caption={Schema der YAML-Evaluationskonfiguration},label={lst:evaluation-config-schema}]
{
    "$schema": "https://json-schema.org/draft/2020-12/schema",
    "$ref": "#/definitions/Configuration",
    "definitions": {
        "Configuration": {
            "type": "object",
            "additionalProperties": false,
            "properties": {
                "defaultEvaluationEndpoint": {
                    "type": "string"
                },
                "maxConcurrent": { "type": "integer" },
                "models": {
                    "type": "array",
                    "items": { "$ref": "#/definitions/Model" }
                },
                "datasets": {
                    "type": "array",
                    "items": { "type": "integer" }
                },
                "seed": { "type": "integer" }
            },
            "required": [
                "defaultEvaluationEndpoint",
                "models"
                "datasets",
            ],
            "title": "Configuration"
        },
        "Model": {
            "type": "object",
            "additionalProperties": false,
            "properties": {
                "label": { "type": "string" },
                "evaluationEndpoint": { "type": "string" },
                "llmProps": { "$ref": "#/definitions/LlmProps" }
            },
            "required": [ "label" ],
            "title": "Model"
        },
        "LlmProps": {
            "type": "object",
            "additionalProperties": false,
            "properties": {
                "baseUrl": {
                    "type": "string",
                    "format": "uri",
                    "qt-uri-protocols": [ "https" ]
                },
                "modelName": { "type": "string" },
                "apiKey": { "type": "string"},
                "timeoutSeconds": { "type": "number" }
            },
            "required": [],
            "title": "LlmProps"
        }
    }
}
\end{lstlisting}

\begin{lstlisting}[caption={Zusammengefasster Logauszug zum Retry-Mechanismus}, label={lst:retry-log}]
2025-10-03T19:11:51.152+02:00  INFO  BpmnExtractor  : Extracting BPMN elements from XML

# 1) Erste Anfrage an das LLM (gekuerzt: Prompt/Headers/Body)
2025-10-03T19:11:51.156+02:00  INFO  LoggingHttpClient  : HTTP POST https://openrouter.ai/api/v1/chat/completions
  model: openai/gpt-oss-20b
  messages: [system: (System-Prompt), user: (User-Prompt mit BpmnElement-Liste und Format-Anweisung)]

# 2) Antwort des LLM mit fehlerhaftem JSON (verkuerzt)
2025-10-03T19:11:56.671+02:00  INFO  LoggingHttpClient  : HTTP 200
assistant:
{
  "elements": [
    { "id": "Activity_09ehuii", "reason": "...", "isRelevant": true },
    { "id": "Activity_1la5hsp", "reason": "...", "isRelevant":  }   <-- fehlender Bool-Wert
    { "id": "Activity_0rfgrlm", "reason": "...", "isRelevant": true }
  ]
}

# 3) Parser-Fehler + Retry-Ankuendigung (gekuerzt)
2025-10-03T19:11:56.691+02:00  WARN  SafetyNet  : Parsing failed. Attempting to fix JSON and retry... (Attempt 1 of 2)
dev.langchain4j.service.output.OutputParsingException:
  Caused by: com.fasterxml.jackson.core.JsonParseException:
  Unexpected character ('}') ... at elements[1].isRelevant

# 4) Zweite Anfrage zum beheben des JSON mit Chat-Verlauf und Fehlermeldung (n-mal wiederholt, bis erfolgreich)
2025-10-03T19:11:56.721+02:00  INFO  LoggingHttpClient  : HTTP POST https://openrouter.ai/api/v1/chat/completions
  messages: [
    system: (System-Prompt),
    user: (User-Prompt mit BpmnElement-Liste und Format-Anweisung),
    assistant: (Fehlerhafte JSON-Antwort),
    system: (Fix-JSON System-Prompt),
    user: (Fehlermeldung)
  ]

# 5) Korrigierte JSON-Antwort des LLM
2025-10-03T19:12:01.519+02:00  INFO  LoggingHttpClient  : HTTP 200
assistant:
{
  "elements": [
    { "id": "Activity_09ehuii", "reason": "...", "isRelevant": true },
    { "id": "Activity_1la5hsp", "reason": "...", "isRelevant": true }, <-- jetzt mit Bool-Wert
    { "id": "Activity_0rfgrlm", "reason": "...", "isRelevant": true }
  ]
}

# 6) Erfolgreiches Parsing und Weiterverarbeitung
2025-10-03T19:12:01.519+02:00  INFO  PromptBpmnAnalyzer : BPMN Analysis Result: elements=[... isRelevant=true ...]
\end{lstlisting}

\backmatter
\nocite{Knappen2009}
\nocite{Mittelbach2005}
\nocite{Schlosser2014}
\nocite{Sturm2012}
\nocite{Voss2010}

\printbibliography

\clearpage
\thispagestyle{empty}

Name: \fullname \hfill Matrikelnummer: \matnr \vspace{2cm}

\minisec{Erklärung}

Ich erkläre, dass ich die Arbeit selbständig verfasst und keine anderen als die angegebenen Quellen und Hilfsmittel verwendet habe.\vspace{2cm}

Ulm, den \dotfill

\hspace{10cm} {\footnotesize \fullname}
\end{document}
